\BeleskeID{09_3-s037}
% element: 09_3-s037:e05
% slide-id: 09_3-s037
Izdvajanjem ulazne kapacitivnosti iz početnog izraza dobijamo
\[\begin{aligned}
t_{p1}&=0.69R_{eq1}(C_{int1}+C_{ext1})\\
&=0.69R_{eq1}C_{in}\left(\frac{C_{ext1}}{C_{in}}+\frac{C_{int1}}{C_{in}}\right).
\end{aligned}\]
Prvi faktor je $\tau_{inv}$. Za ceo lanac isti korak daje
\[t_p=\sum_{i=1}^N0.69R_{eq,i}C_{in,i}
\left(\frac{C_{ext,i}}{C_{in,i}}+\frac{C_{int,i}}{C_{in,i}}\right).
\]
Skaliranje oba tranzistora smanjuje otpornost, a povećava kapacitivnost za isti faktor, pa proizvod ostaje nepromenjen. Time se optimizacija može nastaviti u normalizovanom obliku i dobija se isti rezultat za lanac invertora.
